\section{Bulk statistics and exceptional zeros}\label{sec:complex-statistics}

The transport identities describe a possible defect on the zeta slice.
Before localizing a statistical test around it, we identify what the
established bulk comparisons actually control. Their positivity and
asymptotic formulas concern different observables and different
quantifiers.

\subsection{Coordinates and the original comparison}
Write a nontrivial zeta zero as \(\rho=\sigma+\ii t\), where
\(0<\sigma<1\). To avoid confusing these coordinates with the
collision coefficients \(\beta,\gamma\), we use \(\sigma,t\)
throughout this section. Its corresponding zero of \(H_0\) is
\begin{equation}\label{eq:zeta-H-zero-map}
 z=2t+\ii(1-2\sigma).
\end{equation}
Thus a zero off the critical line is a \emph{nonreal} zero of
\(H_0\), not generally a purely imaginary one. A pure imaginary
zero of \(H_\lambda\), for real \(\lambda\), is impossible:
\(H_\lambda(\ii y)=\int_0^\infty F_\lambda(u)\cosh(yu)\dd u>0\).

Put \(L_T=(\log T)/(2\pi)\), \(B_T=TL_T\), and sum over all
zeros with \(0<t\le T\), counting multiplicity. The ordinate-only
Montgomery sum makes sense without RH:
\begin{equation}\label{eq:F-vertical}
 F_{\rm v}(\alpha,T)=\frac1{B_T}\sum_{\rho,\rho'}
 T^{\ii\alpha(t-t')}\frac4{4+(t-t')^2}.
\end{equation}
Its definition records vertical separation, but discards \(\sigma\).
Montgomery's original restricted-frequency asymptotic assumes RH
\cite[Theorem, equations (1)--(2)]{Montgomery1973}.
The full sine-kernel pair-correlation conjecture is stronger than
that theorem and is not a premise of the Rodgers--Tao proof.
Under \(\Lambda<0\), their dynamics forces almost all gaps at
\(\lambda=0\) to be asymptotic to their mean; their final
contradiction uses a positive proportion of smaller gaps
\cite[Section~9, equation (93)]{RodgersTao2020}. The proof there
uses the gap theorem of Conrey, Ghosh, Goldston, Gonek and
Heath-Brown, rather than assuming the full statistical conjecture.

\subsection{An unconditional complex-difference statistic}
Baluyot, Goldston, Suriajaya and Turnage-Butterbaugh instead retain
the complete difference of zeros. Define
\begin{equation}\label{eq:F-complex}
 F_{\rm c}(\alpha,T)=\frac1{B_T}\sum_{\rho,\rho'}
 T^{\alpha(\rho-\rho')}W(\rho-\rho'),\qquad
 W(u)=\frac4{4-u^2}.
\end{equation}
There is no pole in this sum since \(|\Re(\rho-\rho')|<1\).
Under RH, \(F_{\rm c}=F_{\rm v}\). Without RH they are different:
\(T^{\alpha(\rho-\rho')}\) includes the real factor
\(T^{\alpha(\sigma-\sigma')}\), and the kernel may be complex.

\begin{theorem}[Unconditional Montgomery theorem, imported]
\label{thm:BGST}
For the zeta zeros, \(F_{\rm c}\) is real, even in \(\alpha\),
and nonnegative. As \(T\to\infty\), uniformly for
\(0\le\alpha\le1\),
\begin{equation}\label{eq:BGST}
 F_{\rm c}(\alpha,T)
 =T^{-2\alpha}(\log T+O(1))+\alpha
                        +O((\log T)^{-1/2}).
\end{equation}
No RH assumption is required.
\end{theorem}
This is \cite[Theorem~1, (1.2)--(1.4)]{BGST2024}, in the
pinned arXiv version~1. Its arithmetic asymptotic is an imported
result. The following direct calculation explains why positivity
survives the presence of nonreal zeros.

\begin{proposition}[A positive integral for the finite sum]
\label{prop:complex-Gram}
For any finite multiset \(\mathcal Z\) in \(0<\Re\rho<1\)
invariant under \(\rho\mapsto1-\overline\rho\), and \(x>0\),
\begin{equation}\label{eq:complex-Gram}
 \sum_{\rho,\rho'\in\mathcal Z}x^{\rho-\rho'}W(\rho-\rho')
 =\int_\R\ee^{-2|v|}
 \left|\sum_{\rho\in\mathcal Z}
       \ee^{(\rho-1/2)(\log x+v)}\right|^2\dd v\ge0.
\end{equation}
\end{proposition}
\begin{proof}
For \(|\Re u|<2\), direct integration on the two half-lines gives
\(W(u)=\int_\R\ee^{-2|v|}\ee^{uv}\dd v\).
Relabel \(\rho'\) by \(1-\overline\rho'\). The summand becomes
\(x^{\rho+\overline\rho'-1}W(\rho+\overline\rho'-1)\).
Insert the integral and expand the absolute square. Finite
summation is legitimate; the largest real exponent in the
integrand has absolute value less than one, so the integral
converges. Swapping the two indices in the original expression
also proves invariance under \(x\mapsto1/x\).
\end{proof}

The published proof has an equivalent positive integral with
Cauchy kernels \cite[Lemma~3, (2.6)]{BGST2024}.
Thus the obstacle is not simply loss of positivity of the
statistic. Its summands, however, are not nonnegative pair
counts, and this positivity does not imply that each zero is
on the critical line. Analytic test functions at complex zero
coordinates also occur in the earlier higher-correlation theorem
of Rudnick and Sarnak \cite[Theorem~1.1]{RudnickSarnak1996};
its restricted Fourier support and entire smoothing must be
retained. It is not a theorem for arbitrary planar pair tests.

\subsection{What ordinate statistics can detect}
The functional equation pairs \(\sigma+\ii t\) with
\(1-\sigma+\ii t\). Consequently, even a statistic that discards
\(\sigma\) can detect the coincidence of their ordinates.
For each distinct ordinate \(t\), let \(m_t\) be the total
multiplicity of zeros on that horizontal line, and set
\begin{equation}\label{eq:horizontal-excess}
 N(T)=\sum_t m_t,\qquad D(T)=\sum_t m_t(m_t-1).
\end{equation}
The pair sum including all equal-ordinate terms is \(N(T)+D(T)\).
This is the horizontal-multiplicity observable studied by
Goldston, Lee, Schettler and Suriajaya
\cite[Sections~7--8]{GLSS2026}.

\begin{proposition}[Coincidence bounds]\label{prop:horizontal}
Let \(N_{\rm off}(T)\) count zeros off the critical line, with
multiplicity, and let \(N_{\rm sc}(T)\) count simple critical-line
zeros. Then
\begin{equation}\label{eq:horizontal-bounds}
 N_{\rm off}(T)\le D(T),\qquad
 N_{\rm sc}(T)\ge N(T)-D(T).
\end{equation}
In particular \(D(T)=o(N(T))\) implies that a proportion tending
to one of the zeros is simple and on the critical line.
\end{proposition}
\begin{proof}
Every off-line zero has a distinct reflected partner at the
same ordinate, so lies on a line with \(m_t\ge2\). Such a line
contributes \(m_t(m_t-1)\ge m_t\) to \(D(T)\). A line with
\(m_t=1\) contains one simple zero, which the reflection forces
onto the critical line. The number of copies on lines with
\(m_t\ge2\) is at most \(D(T)\), proving both bounds.
\end{proof}

For precision, the ordinate conjecture used in \cite{GLSS2026}
asserts
\begin{equation}\label{eq:ordinate-PCC}
 \frac1{B_T}\sum_{\substack{0<t,t'\le T\\0<(t'-t)L_T\le r}}1
 \longrightarrow\int_0^r
       \left(1-\left(\frac{\sin\pi u}{\pi u}\right)^2\right)\dd u,
\end{equation}
uniformly for \(r\) in compact subintervals of \((0,\infty)\),
with all zeros counted with multiplicity and equal ordinates
excluded from this particular sum. Their Theorem~1 proves the
density-one conclusion from this conjecture without assuming RH. This remains conditional
on that conjecture; it does not follow from
Theorem~\ref{thm:BGST}. Even the resulting density-one conclusion
allows a finite or zero-density exceptional set, so is weaker
than RH. The elementary implication \eqref{eq:horizontal-bounds}
is useful precisely because it exposes this quantifier.

\subsection{Finite exceptions are invisible to the leading asymptotic}
The limitation can be proved directly for both sums, rather than
inferred only from the language of density.

\begin{proposition}[Stability under finitely many zero changes]
\label{prop:finite-defects}
Let two locally finite multisets in \(0<\sigma<1\), \(t>0\),
differ by finitely many fixed zeros with multiplicity. Suppose
their ordinate counts in \([u,u+1]\) are \(O(\log(2+u))\).
Define \(F_{\rm v},F_{\rm c}\) by \eqref{eq:F-vertical} and
\eqref{eq:F-complex}, with the common normalizer \(B_T\).
If all changed zeros have \(\eta\le\sigma\le1-\eta\) for
some \(\eta>0\), then, uniformly for \(0\le\alpha\le1\),
\begin{equation}\label{eq:finite-defect-errors}
 |\Delta F_{\rm v}|=O((T\log T)^{-1}),\qquad
 |\Delta F_{\rm c}|=O(T^{-\eta}/\log T).
\end{equation}
\end{proposition}
\begin{proof}
After cancellation, each remaining ordered-pair term contains
at least one fixed changed zero \(\rho_0=\sigma_0+\ii t_0\).
The count hypothesis gives
\[
 \sum_\rho\frac1{1+(t-t_0)^2}<\infty:
\]
on distant unit intervals the sum is bounded by a constant times
\(\sum_{n\ge1}\log(2+n)/(1+n^2)\), and there are finitely
many nearby zeros. The vertical phase has modulus one, so
the unnormalized change in \(F_{\rm v}\) is bounded.
For the complex kernel, writing \(d=\sigma-\sigma_0\),
\[
 |4-(d+\ii(t-t_0))^2|
 \ge4-d^2+(t-t_0)^2\ge3+(t-t_0)^2.
\]
Moreover \(|d|\le1-\eta\), whence each exponential factor
has modulus at most \(T^{\alpha(1-\eta)}\).
The finitely many changes therefore alter the unnormalized
complex sum by \(O(T^{1-\eta})\) on the stated frequency
interval. Division by \(B_T\) proves the result.
\end{proof}

The zeta count hypothesis is unconditional
\cite[Lemma~2]{BGST2024}. Symmetric finite changes can be used
if one wishes to preserve the reflection in
Proposition~\ref{prop:complex-Gram}. This proposition concerns
the sensitivity of the observables; it constructs no alternative
zeta function. It shows that their leading limiting formulas,
alone, cannot exclude a finite exceptional set. Frequencies
beyond the proved range, more precise remainders, or localized
observables require their own arithmetic estimates.

\subsection{Statistics that retain the transverse coordinates}
A direct geometric generalization is to keep the transverse
coordinate as a mark. For the zeros with \(T<t_j\le2T\), set
\(a_j=L_Tt_j\), \(b_j=L_T(\sigma_j-1/2)\), and let
\(n_T\) be their number. The finite positive pair measure
\begin{equation}\label{eq:marked-pair}
 \mathcal R_T=\frac1{n_T}\sum_{j\ne k}
             \delta_{(a_j-a_k,b_j,b_k)}
\end{equation}
retains both marks, including whether a pair is symmetric across
the critical line. On RH, testing it against \(f(s,b,b')\)
reduces to the usual ordinate-pair test \(f(s,0,0)\).
Its full finite structure factor is also nonnegative:
\begin{equation}\label{eq:marked-structure}
 \mathcal S_T(k,l)=\frac1{n_T}
 \left|\sum_j\ee^{2\pi\ii(ka_j+lb_j)}\right|^2\ge0.
\end{equation}
This is a two-dimensional Fourier statistic of the zero counting
measure. It is different from the analytic complex-difference
sum \eqref{eq:F-complex}, and no sine-kernel limit for arbitrary
marked tests is asserted. The theorem above does not supply one.

For a finite window of \(H_\lambda\) zeros, one can likewise
use \(\Re z\) and \(\Im z\) as longitudinal and transverse
coordinates. At \(\lambda=0\), \eqref{eq:zeta-H-zero-map}
fixes both scale factors. For positive time, the appropriate
asymptotic real spacing is described next. These finite
definitions need no assumption that all zeros are real.

\subsection{Why the positive-threshold argument is different}
There is a particularly strong constraint on any attempt to
reverse the Rodgers--Tao contradiction.

\begin{theorem}[Positive-time tail, imported]\label{thm:KKL}
For every fixed \(\lambda>0\), all but finitely many zeros
of the canonical \(H_\lambda\) are real and simple. Enumerating
the sufficiently large positive real zeros by \(x_n(\lambda)\),
\begin{equation}\label{eq:positive-clock}
 \frac{(x_{n+1}(\lambda)-x_n(\lambda))\log x_n(\lambda)}{4\pi}
 \longrightarrow1.
\end{equation}
\end{theorem}
These are Ki--Kim--Lee's Theorems~1.3 and~1.4
\cite{KiKimLee2009}. Their normalization is
\(\Xi_\mu(t)=8H_{4\mu}(2t)\); thus \(\mu=\lambda/4\)
and the gap constant in our coordinate is \(4\pi\).
The theorem is for fixed positive time. It gives no uniform
tail threshold as \(\lambda\downarrow0\).

If \(\Lambda>0\), choose \(0<\lambda<\Lambda\).
By definition \(H_\lambda\) has a nonreal zero, but
Theorem~\ref{thm:KKL} confines all its nonreal zeros to a
finite exceptional set. The asymptotic real gaps still satisfy
\eqref{eq:positive-clock}. Thus regular spacing at a positive
time is compatible with either side of the putative threshold;
it is not a contradiction. For a sufficiently high window,
the marked measure \eqref{eq:marked-pair}, adapted to this
fixed slice, contains no transverse marks even if exceptional
nonreal zeros remain below it. Passing from fixed positive
time to \(\lambda=0\) requires a uniform estimate in a joint
time-height limit; compact convergence of the entire functions
alone does not supply it.

The dynamics in Section~\ref{sec:complex-dynamics} identifies the
exterior term needed to transport a particular exception. A statistical
argument must then retain its signal after background interference or
center averaging. We now make those two requirements quantitative.
